\documentclass[11pt,a4paper]{article}

\usepackage[margin=2.2cm]{geometry}
\usepackage{fancyhdr}
\usepackage{listings}
\usepackage{xcolor}
\usepackage{hyperref}
\usepackage{titlesec}
\usepackage{enumitem}
\usepackage{caption}

\setlength{\headheight}{14pt}
\pagestyle{fancy}
\fancyhf{}
\fancyhead[L]{Building a Minimal Command-Line Web Client in C}
\fancyhead[R]{\thepage}
\renewcommand{\headrulewidth}{0.4pt}

\titleformat{\section}{\Large\bfseries\color{blue!70!black}}{}{0em}{}[\titlerule]
\titleformat{\subsection}{\large\bfseries\color{blue!50!black}}{}{0em}{}

% Colors for progressive source listings
\definecolor{codeold}{gray}{0.58}
\definecolor{codebg}{RGB}{248,248,248}
\definecolor{newbg}{RGB}{232,245,233}
\definecolor{newframe}{RGB}{46,125,50}

\lstdefinestyle{basec}{
    language=C,
    basicstyle=\ttfamily\footnotesize,
    breaklines=true,
    breakatwhitespace=false,
    columns=fullflexible,
    keepspaces=true,
    numbers=left,
    numberstyle=\tiny\color{gray},
    numbersep=6pt,
    tabsize=4,
    showstringspaces=false,
    showspaces=false,
    captionpos=b,
    xleftmargin=1.2em,
    framexleftmargin=0.6em,
    aboveskip=0.4em,
    belowskip=0.15em
}

% Prior / unchanged code: light gray so it recedes
\lstdefinestyle{oldcode}{
    style=basec,
    basicstyle=\ttfamily\footnotesize\color{codeold},
    backgroundcolor=\color{codebg},
    keywordstyle=\color{codeold},
    commentstyle=\color{codeold},
    stringstyle=\color{codeold},
    identifierstyle=\color{codeold},
    frame=none,
    rulecolor=\color{gray!30}
}

% Newly added code: full color so it pops
\lstdefinestyle{newcode}{
    style=basec,
    basicstyle=\ttfamily\footnotesize\color{black},
    backgroundcolor=\color{newbg},
    keywordstyle=\color{blue!80!black},
    commentstyle=\color{green!45!black},
    stringstyle=\color{red!70!black},
    identifierstyle=\color{black},
    frame=leftline,
    framerule=2.5pt,
    rulecolor=\color{newframe}
}

% Standalone \"new code\" snippet (not part of full-file collage)
\lstdefinestyle{snippet}{
    style=basec,
    basicstyle=\ttfamily\footnotesize,
    backgroundcolor=\color{newbg},
    keywordstyle=\color{blue!80!black},
    commentstyle=\color{green!45!black},
    stringstyle=\color{red!70!black},
    frame=single,
    rulecolor=\color{newframe},
    aboveskip=0.8em,
    belowskip=0.8em
}

% Final complete file (everything \"new\" / current)
\lstdefinestyle{fullcode}{
    style=basec,
    basicstyle=\ttfamily\footnotesize,
    backgroundcolor=\color{gray!6},
    keywordstyle=\color{blue!80!black},
    commentstyle=\color{green!45!black},
    stringstyle=\color{red!70!black},
    frame=single,
    rulecolor=\color{gray!50},
    aboveskip=0.8em,
    belowskip=0.8em
}

% Continuous line numbers across old/new listing fragments of one file
% (Do not name these \end... --- LaTeX reserves that prefix for environments.)
\newcommand{\FullSourceStart}[1]{%
  \subsection{Full Source at This Stage}%
  \textit{Gray = already present. Green-highlighted = added in this step.}%
  \par\vspace{0.4em}%
  \lstset{name=#1}%
}
\newcommand{\FullSourceStop}{\lstset{name={}}}

\title{\textbf{Building a Minimal Command-Line Internet Client in C}\\[0.4em]
\large A Step-by-Step Tutorial}
\author{S.F.}
\date{July 2026}

\begin{document}

\maketitle

\begin{abstract}
This tutorial guides you through the creation of a compact yet powerful
command-line tool (\texttt{urlget}) capable of fetching web pages over HTTP
and HTTPS, with optional Tor support via SOCKS5. Development is incremental:
each step adds a focused piece of functionality, explains \emph{why} it exists,
and shows the complete source so far---with prior code muted in gray and new
code highlighted---so you always know what changed and where it lives.
\end{abstract}

\section*{Foreword}
\addcontentsline{toc}{section}{Foreword}

High-level libraries are wonderful until you need to know what they are
actually doing. This tutorial takes the opposite path: we talk to the network
with ordinary C, POSIX sockets, and a thin slice of OpenSSL. The point is not
machismo. It is demystification.

\textbf{Almost bare-metal networking is not hard.} Creating a socket,
connecting it, writing bytes, and reading bytes is a short, repeatable recipe.
Once those four moves are clear, you already understand the bulk of what every
network client does under the hood. Making the same path meaningfully faster
in user space usually means dropping to assembly (or accepting that the kernel,
the NIC, and physics own the rest of the budget). Clarity first; micro-optimizations
later, and only if measurement demands them.

\textbf{After the building blocks, the pattern barely changes.} Resolve an
address, open a socket, connect, send, receive, close. That loop is the same
whether the peer speaks HTTP, SMTP, a binary game protocol, or something you
invented this afternoon. What changes in the later stages is almost always the
\emph{payload}: is the protocol text or binary, how are messages framed, and
how do you parse them safely? HTTP headers, SOCKS5 handshakes, and TLS records
look different on the wire, but they all ride the same socket primitives.

\texttt{urlget} is a small vehicle for that idea---a mini command-line client
you can read end to end. If you finish it and think ``that was mostly sockets
plus parsing,'' the foreword did its job.

\hfill--- S.F.

\tableofcontents
\newpage

\section{Introduction and Goals}

We will build a tool called \texttt{urlget}: a small C program that behaves
like a minimal subset of \texttt{curl}. By the end of this tutorial you will
have a single-file client that can:

\begin{itemize}[leftmargin=1.4em]
    \item Fetch content from \texttt{http://} and \texttt{https://} URLs
    \item Optionally route traffic through Tor via SOCKS5
    \item Follow HTTP redirects (with a configurable limit)
    \item Send custom headers, POST bodies, and basic authentication
    \item Write the response body to a file instead of stdout
    \item Compile clean release builds with debug tracing disabled via
          \texttt{-DNDEBUG}
\end{itemize}

\subsection{Design Principles}

\begin{enumerate}[leftmargin=1.4em]
    \item \textbf{Incremental.} Each step produces a coherent program. You can
          stop after any step and still understand what you have.
    \item \textbf{Readable over clever.} Prefer clear control flow and small
          helpers over dense one-liners.
    \item \textbf{POSIX + OpenSSL.} We use standard sockets for TCP and
          OpenSSL 3 for TLS. No third-party HTTP libraries.
    \item \textbf{Explain the \emph{why}.} Every new block is introduced with
          the problem it solves, not only the syntax.
\end{enumerate}

\subsection{What You Need}

\begin{itemize}[leftmargin=1.4em]
    \item A C compiler (\texttt{gcc} or \texttt{clang})
    \item OpenSSL development headers (for the HTTPS step onward), e.g.\
          \texttt{libssl-dev} on Debian/Ubuntu
    \item Optional: a local Tor daemon listening on \texttt{127.0.0.1:9050}
          if you want to try the SOCKS5 path
\end{itemize}

\subsection{How to Read This Tutorial}

Each development step has the same shape:

\begin{enumerate}[leftmargin=1.4em]
    \item A short motivation: what is missing and why we add it now
    \item Focused snippets of the \textbf{new} code (green background)
    \item The \textbf{full source so far}: unchanged lines in light gray,
          newly added lines in green so they stand out in context
\end{enumerate}

The final step assembles the complete \texttt{urlget.c} matching a real,
usable client.

\section{Step 1: Basic Project Skeleton}

Every non-trivial program starts as something almost trivial. Before sockets,
TLS, or argument parsing, we need a file that compiles, accepts a URL on the
command line, and proves the tool is wired up.

\subsection{Why This Step Exists}

If \texttt{main} cannot see \texttt{argv}, nothing later matters. A skeleton
also gives you a compile command you will reuse for the rest of the tutorial:

\begin{lstlisting}[style=snippet,language=bash]
gcc -Wall -Wextra -O2 -o urlget urlget.c
./urlget http://example.com
\end{lstlisting}

\subsection{What We Add}

\begin{itemize}[leftmargin=1.4em]
    \item Standard headers for I/O and \texttt{EXIT\_FAILURE}
    \item A check that at least one argument (the URL) was provided
    \item A confirmation print so you can see the argument was received
\end{itemize}

\subsection{New Code}

\begin{lstlisting}[style=snippet,language=C]
#include <stdio.h>
#include <stdlib.h>

int main(int argc, char *argv[])
{
    if (argc < 2) {
        fprintf(stderr, "Usage: %s <URL>\n", argv[0]);
        return EXIT_FAILURE;
    }
    printf("Requested: %s\n", argv[1]);
    return 0;
}
\end{lstlisting}

Error messages go to \texttt{stderr}; successful chatter goes to
\texttt{stdout}. That split stays useful once the program starts printing
HTTP bodies to stdout.

\FullSourceStart{urlgetstep01}
\begin{lstlisting}[style=newcode,language=C,firstnumber=1]
/* urlget.c -- Step 1: basic skeleton */
#include <stdio.h>
#include <stdlib.h>

int main(int argc, char *argv[])
{
    if (argc < 2) {
        fprintf(stderr, "Usage: %s <URL>\n", argv[0]);
        return EXIT_FAILURE;
    }
    printf("Requested: %s\n", argv[1]);
    return 0;
}
\end{lstlisting}
\FullSourceStop

At this stage everything is new, so the whole file is highlighted. From the
next step onward, only the deltas stay green.

\section{Step 2: A Clean \texttt{usage()} Function}

Hard-coding the usage string inside \texttt{main} works once. As soon as we
add options, the same message must be printed from many places (missing
argument, missing option value, unknown flags). Duplicating the string is how
usage text drifts out of date.

\subsection{Why This Step Exists}

\begin{itemize}[leftmargin=1.4em]
    \item \textbf{Single source of truth} for how to invoke the program
    \item \textbf{Immediate exit} on misuse---no half-initialized state
    \item \texttt{static inline} keeps a tiny helper local to this translation
          unit without the overhead of a separate \texttt{.c} file
\end{itemize}

\subsection{What Changes}

We extract the error path into \texttt{usage()}, then call it when the URL is
missing. The temporary \texttt{printf} confirmation stays for now so you can
still smoke-test the binary; it will disappear once we actually fetch pages.

\subsection{New Code}

\begin{lstlisting}[style=snippet,language=C]
static inline void usage(const char *progname)
{
    fprintf(stderr, "Usage: %s <URL>\n", progname);
    exit(EXIT_FAILURE);
}
\end{lstlisting}

And in \texttt{main}:

\begin{lstlisting}[style=snippet,language=C]
if (argc < 2)
    usage(argv[0]);
\end{lstlisting}

Passing \texttt{argv[0]} means the message shows whatever name the user
typed (\texttt{./urlget}, \texttt{urlget}, a full path, \ldots).

\FullSourceStart{urlgetstep02}
\begin{lstlisting}[style=oldcode,language=C,firstnumber=1]
/* urlget.c -- Step 2: usage() helper */
#include <stdio.h>
#include <stdlib.h>
\end{lstlisting}
\begin{lstlisting}[style=newcode,language=C,firstnumber=last]
static inline void usage(const char *progname)
{
    fprintf(stderr, "Usage: %s <URL>\n", progname);
    exit(EXIT_FAILURE);
}
\end{lstlisting}
\begin{lstlisting}[style=oldcode,language=C,firstnumber=last]
int main(int argc, char *argv[])
{
\end{lstlisting}
\begin{lstlisting}[style=newcode,language=C,firstnumber=last]
    if (argc < 2)
        usage(argv[0]);
\end{lstlisting}
\begin{lstlisting}[style=oldcode,language=C,firstnumber=last]
    printf("Requested: %s\n", argv[1]);
    return 0;
}
\end{lstlisting}
\FullSourceStop

\section{Step 3: Basic Argument Parsing}

Real clients accept flags before the URL. We introduce a straightforward
\texttt{for} loop over \texttt{argv} rather than \texttt{getopt\_long}, so
every branch is visible and easy to extend in later steps.

\subsection{Why a Manual Loop?}

\begin{itemize}[leftmargin=1.4em]
    \item No extra dependencies or platform quirks
    \item Optional values (e.g.\ \texttt{-T} with or without a proxy) are
          easier to express than with classic \texttt{getopt}
    \item Students see exactly when the index advances past an option argument
\end{itemize}

\subsection{State We Introduce}

\begin{itemize}[leftmargin=1.4em]
    \item \texttt{url\_arg} --- the single positional URL (required)
    \item \texttt{user\_agent} --- default browser-like string, overridable
          with \texttt{-A} / \texttt{--user-agent}
    \item Placeholders for Tor and redirect policy (wired up in later steps)
\end{itemize}

We also expand \texttt{usage()} so it documents the options we claim to
support. Empty stubs (\texttt{use\_tor}, \texttt{follow\_redirects}, \ldots)
document the roadmap without implementing behavior yet.

\subsection{New Code --- Expanded Usage}

\begin{lstlisting}[style=snippet,language=C]
static inline void usage(const char *progname)
{
    fprintf(stderr,
        "Usage: %s [options] <URL>\n\n"
        "Options:\n"
        "  -A, --user-agent STR     Set User-Agent\n"
        "  -T, --tor [proxy]        Use Tor (default 127.0.0.1:9050)\n"
        "  -H, --header STR         Add custom header (repeatable)\n"
        "  -d, --data DATA          POST data (@file supported)\n"
        "  -o, --output FILE        Write output to file\n"
        "  -u, --user USER:PASS     Basic authentication\n"
        "  -I, --head               HEAD request\n"
        "      --no-location        Do not follow redirects\n"
        "      --max-redirects N    Max redirects (default: 10)\n",
        progname);
    exit(EXIT_FAILURE);
}
\end{lstlisting}

\subsection{New Code --- Parsing Loop (core options)}

For this step we fully implement \texttt{-A} and recognition of the URL; Tor
and the other flags are parsed enough to set variables we will use later.

\begin{lstlisting}[style=snippet,language=C]
const char *user_agent =
    "Mozilla/5.0 (X11; Linux x86_64; rv:128.0) "
    "Gecko/20100101 Firefox/128.0";
const char *url_arg = NULL;

for (int i = 1; i < argc; i++) {
    if (strcmp(argv[i], "-A") == 0 ||
        strcmp(argv[i], "--user-agent") == 0) {
        if (i + 1 < argc) user_agent = argv[++i];
        else usage(argv[0]);
    }
    else if (url_arg == NULL) {
        url_arg = argv[i];
    }
}

if (url_arg == NULL)
    usage(argv[0]);
\end{lstlisting}

\textbf{Rule of thumb:} option switches that take a value always
\texttt{++i} after validating \texttt{i + 1 < argc}. Anything that is not a
recognized option becomes the URL---first one wins.

\FullSourceStart{urlgetstep03}
\begin{lstlisting}[style=oldcode,language=C,firstnumber=1]
/* urlget.c -- Step 3: argument parsing */
#include <stdio.h>
#include <stdlib.h>
\end{lstlisting}
\begin{lstlisting}[style=newcode,language=C,firstnumber=last]
#include <string.h>
#include <stdbool.h>
\end{lstlisting}
\begin{lstlisting}[style=oldcode,language=C,firstnumber=last]

\end{lstlisting}
\begin{lstlisting}[style=newcode,language=C,firstnumber=last]
static inline void usage(const char *progname)
{
    fprintf(stderr,
        "Usage: %s [options] <URL>\n\n"
        "Options:\n"
        "  -A, --user-agent STR     Set User-Agent\n"
        "  -T, --tor [proxy]        Use Tor (default 127.0.0.1:9050)\n"
        "  -H, --header STR         Add custom header (repeatable)\n"
        "  -d, --data DATA          POST data (@file supported)\n"
        "  -o, --output FILE        Write output to file\n"
        "  -u, --user USER:PASS     Basic authentication\n"
        "  -I, --head               HEAD request\n"
        "      --no-location        Do not follow redirects\n"
        "      --max-redirects N    Max redirects (default: 10)\n",
        progname);
    exit(EXIT_FAILURE);
}
\end{lstlisting}
\begin{lstlisting}[style=oldcode,language=C,firstnumber=last]
int main(int argc, char *argv[])
{
\end{lstlisting}
\begin{lstlisting}[style=newcode,language=C,firstnumber=last]
    const char *user_agent =
        "Mozilla/5.0 (X11; Linux x86_64; rv:128.0) "
        "Gecko/20100101 Firefox/128.0";
    const char *default_tor = "127.0.0.1:9050";

    const char *url_arg = NULL;
    const char *tor_proxy = NULL;
    bool use_tor = false;
    bool follow_redirects = true;
    int max_redirects = 10;
    bool head_request = false;
    char extra_headers[4096] = {0};
    char post_data[8192] = {0};
    char output_file[256] = {0};
    char auth_header[256] = {0};

    for (int i = 1; i < argc; i++) {
        if (strcmp(argv[i], "-A") == 0 ||
            strcmp(argv[i], "--user-agent") == 0) {
            if (i + 1 < argc) user_agent = argv[++i];
            else usage(argv[0]);
        }
        else if (strcmp(argv[i], "-T") == 0 ||
                 strcmp(argv[i], "--tor") == 0) {
            use_tor = true;
            if (i + 1 < argc && argv[i + 1][0] != '-' &&
                strchr(argv[i + 1], ':')) {
                tor_proxy = argv[i + 1];
                i++;
            }
        }
        else if (strcmp(argv[i], "-H") == 0 ||
                 strcmp(argv[i], "--header") == 0) {
            if (i + 1 < argc) {
                strncat(extra_headers, argv[++i],
                        sizeof(extra_headers) -
                        strlen(extra_headers) - 3);
                strcat(extra_headers, "\r\n");
            } else usage(argv[0]);
        }
        else if (strcmp(argv[i], "-d") == 0 ||
                 strcmp(argv[i], "--data") == 0) {
            if (i + 1 < argc) {
                const char *val = argv[++i];
                if (val[0] == '@') {
                    FILE *f = fopen(val + 1, "r");
                    if (f) {
                        fread(post_data, 1,
                              sizeof(post_data) - 1, f);
                        fclose(f);
                    }
                } else {
                    strncpy(post_data, val,
                            sizeof(post_data) - 1);
                }
            } else usage(argv[0]);
        }
        else if (strcmp(argv[i], "-o") == 0 ||
                 strcmp(argv[i], "--output") == 0) {
            if (i + 1 < argc)
                strncpy(output_file, argv[++i],
                        sizeof(output_file) - 1);
            else usage(argv[0]);
        }
        else if (strcmp(argv[i], "-u") == 0 ||
                 strcmp(argv[i], "--user") == 0) {
            if (i + 1 < argc)
                snprintf(auth_header, sizeof(auth_header),
                    "Authorization: Basic %s\r\n", argv[++i]);
            else usage(argv[0]);
        }
        else if (strcmp(argv[i], "-I") == 0 ||
                 strcmp(argv[i], "--head") == 0) {
            head_request = true;
        }
        else if (strcmp(argv[i], "--no-location") == 0) {
            follow_redirects = false;
        }
        else if (strcmp(argv[i], "--max-redirects") == 0 &&
                 i + 1 < argc) {
            max_redirects = atoi(argv[++i]);
        }
        else if (url_arg == NULL) {
            url_arg = argv[i];
        }
    }

    if (url_arg == NULL)
        usage(argv[0]);
    if (use_tor && tor_proxy == NULL)
        tor_proxy = default_tor;

    printf("URL: %s\n", url_arg);
    printf("User-Agent: %s\n", user_agent);
    (void)follow_redirects; (void)max_redirects;
    (void)head_request; (void)extra_headers;
    (void)post_data; (void)output_file; (void)auth_header;
\end{lstlisting}
\begin{lstlisting}[style=oldcode,language=C,firstnumber=last]
    return 0;
}
\end{lstlisting}
\FullSourceStop

The \texttt{(void)} casts silence unused-variable warnings until later steps
consume those fields. Parsing early means step~9 will only \emph{use} the
flags, not reinvent the CLI.

\section{Step 4: URL Parsing (\texttt{parse\_url})}

Sockets need a host and port. HTTP needs a path. The user gives us one string
like \texttt{https://example.com:8443/index.html?q=1}. \texttt{parse\_url}
turns that string into a small structure we can trust for the rest of the
program.

\subsection{Why a Dedicated Struct?}

\begin{lstlisting}[style=snippet,language=C]
typedef struct {
    bool https;
    char host[256];
    uint16_t port;
    char path[512];
} url_t;
\end{lstlisting}

Bundling fields avoids parallel arrays of arguments and makes later redirect
handling easy: parse a new URL into a fresh \texttt{url\_t} and continue.

\subsection{What the Parser Must Handle}

\begin{itemize}[leftmargin=1.4em]
    \item Optional scheme: \texttt{http://} (port~80) or \texttt{https://}
          (port~443)
    \item Host with optional \texttt{:port}
    \item IPv6 literals in brackets: \texttt{[::1]}:port
    \item Path defaulting to \texttt{"/"} when absent
    \item Query-only forms such as \texttt{http://h?x=1} $\rightarrow$ path
          \texttt{/?x=1}
\end{itemize}

\subsection{Algorithm (high level)}

\begin{enumerate}[leftmargin=1.4em]
    \item Copy the argument into a mutable buffer (we will insert
          \texttt{'\textbackslash0'} while splitting).
    \item Detect and strip the scheme; set \texttt{https} and the default port.
    \item Find the first \texttt{/} or \texttt{?} --- that starts the path
          (or query). Everything before it is \texttt{host[:port]}.
    \item Split host and port (special-case \texttt{[ipv6]}).
    \item Reject oversized components rather than overflowing fixed buffers.
\end{enumerate}

\subsection{New Code --- Core of \texttt{parse\_url}}

\begin{lstlisting}[style=snippet,language=C]
static inline int parse_url(const char *arg, url_t *u)
{
    char buf[1024];
    if (strlen(arg) >= sizeof(buf)) return -1;
    strcpy(buf, arg);

    char *p = buf;
    u->https = false;

    if (strncmp(p, "https://", 8) == 0) {
        u->https = true; p += 8;
    } else if (strncmp(p, "http://", 7) == 0) {
        p += 7;
    }

    char *sep = strpbrk(p, "/?");
    char *hostport = p;

    if (sep) {
        if (*sep == '?') {
            /* "host?query" -> path "/?query" */
            size_t ql = strlen(sep);
            if (ql + 2 > sizeof(u->path)) return -1;
            u->path[0] = '/';
            memcpy(u->path + 1, sep, ql + 1);
        } else {
            if (strlen(sep) >= sizeof(u->path)) return -1;
            strcpy(u->path, sep);
        }
        *sep = '\0';
    } else {
        strcpy(u->path, "/");
    }

    /* host / [ipv6] / port split ... */
    uint16_t def = u->https ? 443 : 80;
    u->port = def;
    /* ... see full source below for host/port details ... */
    return 0;
}
\end{lstlisting}

Call it from \texttt{main} after argument parsing succeeds:

\begin{lstlisting}[style=snippet,language=C]
url_t u = {0};
if (parse_url(url_arg, &u) != 0) {
    fprintf(stderr, "Failed to parse URL\n");
    return EXIT_FAILURE;
}
printf("https=%d host=%s port=%u path=%s\n",
       u.https, u.host, u.port, u.path);
\end{lstlisting}

\FullSourceStart{urlgetstep04}
\begin{lstlisting}[style=oldcode,language=C,firstnumber=1]
/* urlget.c -- Step 4: URL parsing */
#include <stdio.h>
#include <stdlib.h>
#include <string.h>
#include <stdbool.h>
\end{lstlisting}
\begin{lstlisting}[style=newcode,language=C,firstnumber=last]
#include <stdint.h>
\end{lstlisting}
\begin{lstlisting}[style=oldcode,language=C,firstnumber=last]

\end{lstlisting}
\begin{lstlisting}[style=newcode,language=C,firstnumber=last]
typedef struct {
    bool https;
    char host[256];
    uint16_t port;
    char path[512];
} url_t;

static inline int parse_url(const char *arg, url_t *u)
{
    char buf[1024];
    if (strlen(arg) >= sizeof(buf)) return -1;
    strcpy(buf, arg);

    char *p = buf;
    u->https = false;

    if (strncmp(p, "https://", 8) == 0) {
        u->https = true; p += 8;
    } else if (strncmp(p, "http://", 7) == 0) {
        p += 7;
    }

    char *sep = strpbrk(p, "/?");
    char *hostport = p;

    if (sep) {
        if (*sep == '?') {
            size_t ql = strlen(sep);
            if (ql + 2 > sizeof(u->path)) return -1;
            u->path[0] = '/';
            memcpy(u->path + 1, sep, ql + 1);
        } else {
            if (strlen(sep) >= sizeof(u->path)) return -1;
            strcpy(u->path, sep);
        }
        *sep = '\0';
    } else {
        strcpy(u->path, "/");
    }

    char *host = u->host;
    size_t hmax = sizeof(u->host) - 1;
    uint16_t def = u->https ? 443 : 80;
    u->port = def;

    if (*hostport == '[') {
        char *cl = strchr(hostport, ']');
        if (!cl) return -1;
        size_t hl = cl - (hostport + 1);
        if (hl >= hmax) return -1;
        memcpy(host, hostport + 1, hl);
        host[hl] = '\0';
        char *aft = cl + 1;
        if (*aft == ':')
            u->port = (uint16_t)strtol(aft + 1, NULL, 10);
    } else {
        char *col = strrchr(hostport, ':');
        if (col) {
            size_t hl = col - hostport;
            if (hl >= hmax) return -1;
            memcpy(host, hostport, hl);
            host[hl] = '\0';
            u->port = (uint16_t)strtol(col + 1, NULL, 10);
        } else {
            if (strlen(hostport) >= hmax) return -1;
            strcpy(host, hostport);
        }
    }

    if (u->port == 0) u->port = def;
    return 0;
}
\end{lstlisting}
\begin{lstlisting}[style=oldcode,language=C,firstnumber=last]
static inline void usage(const char *progname)
{
    fprintf(stderr,
        "Usage: %s [options] <URL>\n\n"
        "Options:\n"
        "  -A, --user-agent STR     Set User-Agent\n"
        "  -T, --tor [proxy]        Use Tor (default 127.0.0.1:9050)\n"
        "  -H, --header STR         Add custom header (repeatable)\n"
        "  -d, --data DATA          POST data (@file supported)\n"
        "  -o, --output FILE        Write output to file\n"
        "  -u, --user USER:PASS     Basic authentication\n"
        "  -I, --head               HEAD request\n"
        "      --no-location        Do not follow redirects\n"
        "      --max-redirects N    Max redirects (default: 10)\n",
        progname);
    exit(EXIT_FAILURE);
}

int main(int argc, char *argv[])
{
    const char *user_agent =
        "Mozilla/5.0 (X11; Linux x86_64; rv:128.0) "
        "Gecko/20100101 Firefox/128.0";
    const char *default_tor = "127.0.0.1:9050";

    const char *url_arg = NULL;
    const char *tor_proxy = NULL;
    bool use_tor = false;
    bool follow_redirects = true;
    int max_redirects = 10;
    bool head_request = false;
    char extra_headers[4096] = {0};
    char post_data[8192] = {0};
    char output_file[256] = {0};
    char auth_header[256] = {0};

    for (int i = 1; i < argc; i++) {
        if (strcmp(argv[i], "-A") == 0 ||
            strcmp(argv[i], "--user-agent") == 0) {
            if (i + 1 < argc) user_agent = argv[++i];
            else usage(argv[0]);
        }
        else if (strcmp(argv[i], "-T") == 0 ||
                 strcmp(argv[i], "--tor") == 0) {
            use_tor = true;
            if (i + 1 < argc && argv[i + 1][0] != '-' &&
                strchr(argv[i + 1], ':')) {
                tor_proxy = argv[i + 1];
                i++;
            }
        }
        else if (strcmp(argv[i], "-H") == 0 ||
                 strcmp(argv[i], "--header") == 0) {
            if (i + 1 < argc) {
                strncat(extra_headers, argv[++i],
                        sizeof(extra_headers) -
                        strlen(extra_headers) - 3);
                strcat(extra_headers, "\r\n");
            } else usage(argv[0]);
        }
        else if (strcmp(argv[i], "-d") == 0 ||
                 strcmp(argv[i], "--data") == 0) {
            if (i + 1 < argc) {
                const char *val = argv[++i];
                if (val[0] == '@') {
                    FILE *f = fopen(val + 1, "r");
                    if (f) {
                        fread(post_data, 1,
                              sizeof(post_data) - 1, f);
                        fclose(f);
                    }
                } else {
                    strncpy(post_data, val,
                            sizeof(post_data) - 1);
                }
            } else usage(argv[0]);
        }
        else if (strcmp(argv[i], "-o") == 0 ||
                 strcmp(argv[i], "--output") == 0) {
            if (i + 1 < argc)
                strncpy(output_file, argv[++i],
                        sizeof(output_file) - 1);
            else usage(argv[0]);
        }
        else if (strcmp(argv[i], "-u") == 0 ||
                 strcmp(argv[i], "--user") == 0) {
            if (i + 1 < argc)
                snprintf(auth_header, sizeof(auth_header),
                    "Authorization: Basic %s\r\n", argv[++i]);
            else usage(argv[0]);
        }
        else if (strcmp(argv[i], "-I") == 0 ||
                 strcmp(argv[i], "--head") == 0) {
            head_request = true;
        }
        else if (strcmp(argv[i], "--no-location") == 0) {
            follow_redirects = false;
        }
        else if (strcmp(argv[i], "--max-redirects") == 0 &&
                 i + 1 < argc) {
            max_redirects = atoi(argv[++i]);
        }
        else if (url_arg == NULL) {
            url_arg = argv[i];
        }
    }

    if (url_arg == NULL)
        usage(argv[0]);
    if (use_tor && tor_proxy == NULL)
        tor_proxy = default_tor;
\end{lstlisting}
\begin{lstlisting}[style=newcode,language=C,firstnumber=last]
    url_t u = {0};
    if (parse_url(url_arg, &u) != 0) {
        fprintf(stderr, "Failed to parse URL\n");
        return EXIT_FAILURE;
    }
    printf("https=%d host=%s port=%u path=%s\n",
           u.https, u.host, u.port, u.path);
    (void)user_agent; (void)use_tor; (void)tor_proxy;
    (void)follow_redirects; (void)max_redirects;
    (void)head_request; (void)extra_headers;
    (void)post_data; (void)output_file; (void)auth_header;
\end{lstlisting}
\begin{lstlisting}[style=oldcode,language=C,firstnumber=last]
    return 0;
}
\end{lstlisting}
\FullSourceStop

Try a few inputs:

\begin{lstlisting}[style=snippet,language=bash]
./urlget http://example.com
./urlget https://example.com:8443/api?x=1
./urlget http://[::1]:8080/status
\end{lstlisting}

\section{Step 5: TCP Connection + Plain HTTP}

With a parsed host, port, and path we can finally speak HTTP. This step stays
on cleartext \texttt{http://} only; TLS arrives next. The flow is the classic
Berkeley sockets path: resolve $\rightarrow$ socket $\rightarrow$ connect
$\rightarrow$ write request $\rightarrow$ read response.

\subsection{Why These APIs?}

\begin{itemize}[leftmargin=1.4em]
    \item \texttt{getaddrinfo} replaces obsolete \texttt{gethostbyname} and
          supports IPv4 and IPv6 transparently (\texttt{AF\_UNSPEC}).
    \item Trying every result in the addrinfo list handles multi-homed hosts.
    \item HTTP/1.1 with \texttt{Connection: close} keeps the client simple:
          we read until EOF instead of parsing \texttt{Content-Length} or
          chunked encoding.
\end{itemize}

\subsection{Building the Request}

A minimal GET looks like:

\begin{lstlisting}[style=snippet,language=C]
GET /path HTTP/1.1\r\n
Host: example.com\r\n
User-Agent: ...\r\n
Accept: */*\r\n
Connection: close\r\n
\r\n
\end{lstlisting}

Note the blank line (\texttt{\textbackslash r\textbackslash n} alone) that
terminates the headers. Without it, many servers wait forever for more input.

\subsection{New Code --- Resolve, Connect, Request, Read}

\begin{lstlisting}[style=snippet,language=C]
struct addrinfo hints = {0};
hints.ai_family   = AF_UNSPEC;
hints.ai_socktype = SOCK_STREAM;
char ps[6];
snprintf(ps, sizeof(ps), "%u", u.port);

struct addrinfo *res = NULL;
if (getaddrinfo(u.host, ps, &hints, &res) != 0) {
    perror("getaddrinfo");
    return EXIT_FAILURE;
}

int s = -1;
for (struct addrinfo *a = res; a; a = a->ai_next) {
    s = socket(a->ai_family, a->ai_socktype, a->ai_protocol);
    if (s < 0) continue;
    if (connect(s, a->ai_addr, a->ai_addrlen) == 0) break;
    close(s); s = -1;
}
freeaddrinfo(res);
if (s < 0) { perror("connect"); return EXIT_FAILURE; }

char req[4096];
snprintf(req, sizeof(req),
    "GET %s HTTP/1.1\r\n"
    "Host: %s\r\n"
    "User-Agent: %s\r\n"
    "Accept: */*\r\n"
    "Connection: close\r\n"
    "\r\n",
    u.path, u.host, user_agent);
write(s, req, strlen(req));

char buf[8192];
ssize_t n;
while ((n = read(s, buf, sizeof(buf))) > 0)
    fwrite(buf, 1, n, stdout);
close(s);
\end{lstlisting}

If \texttt{u.https} is set we refuse for now---HTTPS needs OpenSSL next.

\FullSourceStart{urlgetstep05}
\begin{lstlisting}[style=oldcode,language=C,firstnumber=1]
/* urlget.c -- Step 5: TCP + plain HTTP */
#include <stdio.h>
#include <stdlib.h>
#include <string.h>
#include <stdbool.h>
#include <stdint.h>
\end{lstlisting}
\begin{lstlisting}[style=newcode,language=C,firstnumber=last]
#include <unistd.h>
#include <sys/types.h>
#include <sys/socket.h>
#include <netdb.h>
\end{lstlisting}
\begin{lstlisting}[style=oldcode,language=C,firstnumber=last]

typedef struct {
    bool https;
    char host[256];
    uint16_t port;
    char path[512];
} url_t;

static inline int parse_url(const char *arg, url_t *u)
{
    char buf[1024];
    if (strlen(arg) >= sizeof(buf)) return -1;
    strcpy(buf, arg);

    char *p = buf;
    u->https = false;

    if (strncmp(p, "https://", 8) == 0) {
        u->https = true; p += 8;
    } else if (strncmp(p, "http://", 7) == 0) {
        p += 7;
    }

    char *sep = strpbrk(p, "/?");
    char *hostport = p;

    if (sep) {
        if (*sep == '?') {
            size_t ql = strlen(sep);
            if (ql + 2 > sizeof(u->path)) return -1;
            u->path[0] = '/';
            memcpy(u->path + 1, sep, ql + 1);
        } else {
            if (strlen(sep) >= sizeof(u->path)) return -1;
            strcpy(u->path, sep);
        }
        *sep = '\0';
    } else {
        strcpy(u->path, "/");
    }

    char *host = u->host;
    size_t hmax = sizeof(u->host) - 1;
    uint16_t def = u->https ? 443 : 80;
    u->port = def;

    if (*hostport == '[') {
        char *cl = strchr(hostport, ']');
        if (!cl) return -1;
        size_t hl = cl - (hostport + 1);
        if (hl >= hmax) return -1;
        memcpy(host, hostport + 1, hl);
        host[hl] = '\0';
        char *aft = cl + 1;
        if (*aft == ':')
            u->port = (uint16_t)strtol(aft + 1, NULL, 10);
    } else {
        char *col = strrchr(hostport, ':');
        if (col) {
            size_t hl = col - hostport;
            if (hl >= hmax) return -1;
            memcpy(host, hostport, hl);
            host[hl] = '\0';
            u->port = (uint16_t)strtol(col + 1, NULL, 10);
        } else {
            if (strlen(hostport) >= hmax) return -1;
            strcpy(host, hostport);
        }
    }

    if (u->port == 0) u->port = def;
    return 0;
}

static inline void usage(const char *progname)
{
    fprintf(stderr,
        "Usage: %s [options] <URL>\n\n"
        "Options:\n"
        "  -A, --user-agent STR     Set User-Agent\n"
        "  -T, --tor [proxy]        Use Tor (default 127.0.0.1:9050)\n"
        "  -H, --header STR         Add custom header (repeatable)\n"
        "  -d, --data DATA          POST data (@file supported)\n"
        "  -o, --output FILE        Write output to file\n"
        "  -u, --user USER:PASS     Basic authentication\n"
        "  -I, --head               HEAD request\n"
        "      --no-location        Do not follow redirects\n"
        "      --max-redirects N    Max redirects (default: 10)\n",
        progname);
    exit(EXIT_FAILURE);
}

int main(int argc, char *argv[])
{
    const char *user_agent =
        "Mozilla/5.0 (X11; Linux x86_64; rv:128.0) "
        "Gecko/20100101 Firefox/128.0";
    const char *default_tor = "127.0.0.1:9050";

    const char *url_arg = NULL;
    const char *tor_proxy = NULL;
    bool use_tor = false;
    bool follow_redirects = true;
    int max_redirects = 10;
    bool head_request = false;
    char extra_headers[4096] = {0};
    char post_data[8192] = {0};
    char output_file[256] = {0};
    char auth_header[256] = {0};

    for (int i = 1; i < argc; i++) {
        if (strcmp(argv[i], "-A") == 0 ||
            strcmp(argv[i], "--user-agent") == 0) {
            if (i + 1 < argc) user_agent = argv[++i];
            else usage(argv[0]);
        }
        else if (strcmp(argv[i], "-T") == 0 ||
                 strcmp(argv[i], "--tor") == 0) {
            use_tor = true;
            if (i + 1 < argc && argv[i + 1][0] != '-' &&
                strchr(argv[i + 1], ':')) {
                tor_proxy = argv[i + 1];
                i++;
            }
        }
        else if (strcmp(argv[i], "-H") == 0 ||
                 strcmp(argv[i], "--header") == 0) {
            if (i + 1 < argc) {
                strncat(extra_headers, argv[++i],
                        sizeof(extra_headers) -
                        strlen(extra_headers) - 3);
                strcat(extra_headers, "\r\n");
            } else usage(argv[0]);
        }
        else if (strcmp(argv[i], "-d") == 0 ||
                 strcmp(argv[i], "--data") == 0) {
            if (i + 1 < argc) {
                const char *val = argv[++i];
                if (val[0] == '@') {
                    FILE *f = fopen(val + 1, "r");
                    if (f) {
                        fread(post_data, 1,
                              sizeof(post_data) - 1, f);
                        fclose(f);
                    }
                } else {
                    strncpy(post_data, val,
                            sizeof(post_data) - 1);
                }
            } else usage(argv[0]);
        }
        else if (strcmp(argv[i], "-o") == 0 ||
                 strcmp(argv[i], "--output") == 0) {
            if (i + 1 < argc)
                strncpy(output_file, argv[++i],
                        sizeof(output_file) - 1);
            else usage(argv[0]);
        }
        else if (strcmp(argv[i], "-u") == 0 ||
                 strcmp(argv[i], "--user") == 0) {
            if (i + 1 < argc)
                snprintf(auth_header, sizeof(auth_header),
                    "Authorization: Basic %s\r\n", argv[++i]);
            else usage(argv[0]);
        }
        else if (strcmp(argv[i], "-I") == 0 ||
                 strcmp(argv[i], "--head") == 0) {
            head_request = true;
        }
        else if (strcmp(argv[i], "--no-location") == 0) {
            follow_redirects = false;
        }
        else if (strcmp(argv[i], "--max-redirects") == 0 &&
                 i + 1 < argc) {
            max_redirects = atoi(argv[++i]);
        }
        else if (url_arg == NULL) {
            url_arg = argv[i];
        }
    }

    if (url_arg == NULL)
        usage(argv[0]);
    if (use_tor && tor_proxy == NULL)
        tor_proxy = default_tor;

    url_t u = {0};
    if (parse_url(url_arg, &u) != 0) {
        fprintf(stderr, "Failed to parse URL\n");
        return EXIT_FAILURE;
    }
\end{lstlisting}
\begin{lstlisting}[style=newcode,language=C,firstnumber=last]
    if (u.https) {
        fprintf(stderr, "HTTPS not implemented yet\n");
        return EXIT_FAILURE;
    }

    struct addrinfo hints = {0};
    hints.ai_family   = AF_UNSPEC;
    hints.ai_socktype = SOCK_STREAM;
    char ps[6];
    snprintf(ps, sizeof(ps), "%u", u.port);

    struct addrinfo *res = NULL;
    if (getaddrinfo(u.host, ps, &hints, &res) != 0) {
        perror("getaddrinfo");
        return EXIT_FAILURE;
    }

    int s = -1;
    for (struct addrinfo *a = res; a; a = a->ai_next) {
        s = socket(a->ai_family, a->ai_socktype, a->ai_protocol);
        if (s < 0) continue;
        if (connect(s, a->ai_addr, a->ai_addrlen) == 0) break;
        close(s); s = -1;
    }
    freeaddrinfo(res);
    if (s < 0) { perror("connect"); return EXIT_FAILURE; }

    char req[4096];
    snprintf(req, sizeof(req),
        "GET %s HTTP/1.1\r\n"
        "Host: %s\r\n"
        "User-Agent: %s\r\n"
        "Accept: */*\r\n"
        "Connection: close\r\n"
        "\r\n",
        u.path, u.host, user_agent);
    write(s, req, strlen(req));

    char buf[8192];
    ssize_t n;
    while ((n = read(s, buf, sizeof(buf))) > 0)
        fwrite(buf, 1, n, stdout);
    close(s);

    (void)use_tor; (void)tor_proxy;
    (void)follow_redirects; (void)max_redirects;
    (void)head_request; (void)extra_headers;
    (void)post_data; (void)output_file; (void)auth_header;
\end{lstlisting}
\begin{lstlisting}[style=oldcode,language=C,firstnumber=last]
    return 0;
}
\end{lstlisting}
\FullSourceStop

Test with a plain HTTP site (many sites redirect to HTTPS---that is fine; you
will still see the redirect response headers and body):

\begin{lstlisting}[style=snippet,language=bash]
./urlget http://example.com/
\end{lstlisting}

\section{Step 6: Adding HTTPS Support}

Most of the modern web is TLS. We wrap the connected TCP socket with OpenSSL
so the same HTTP request bytes travel encrypted. Certificate verification is
intentionally disabled here for a minimal tutorial client
(\texttt{SSL\_VERIFY\_NONE}); production tools should verify the peer.

\subsection{Why a Helper?}

\texttt{https\_setup} isolates all OpenSSL ceremony:

\begin{enumerate}[leftmargin=1.4em]
    \item Create a client \texttt{SSL\_CTX} with \texttt{TLS\_client\_method()}
    \item Allocate an \texttt{SSL} object and attach the socket fd
    \item Set SNI (\texttt{SSL\_set\_tlsext\_host\_name}) so virtual hosts work
    \item Run the handshake with \texttt{SSL\_connect}
\end{enumerate}

On success it returns both \texttt{ctx} and \texttt{ssl} to the caller for
later cleanup.

\subsection{New Code --- \texttt{https\_setup}}

\begin{lstlisting}[style=snippet,language=C]
static inline bool https_setup(int sock, const char *hostname,
                               SSL_CTX **out_ctx, SSL **out_ssl)
{
    SSL_CTX *ctx = SSL_CTX_new(TLS_client_method());
    if (!ctx) { ERR_print_errors_fp(stderr); return false; }

    SSL_CTX_set_verify(ctx, SSL_VERIFY_NONE, NULL);

    SSL *ssl = SSL_new(ctx);
    if (!ssl) { SSL_CTX_free(ctx); return false; }

    SSL_set_fd(ssl, sock);
    SSL_set_tlsext_host_name(ssl, hostname);

    if (SSL_connect(ssl) <= 0) {
        ERR_print_errors_fp(stderr);
        SSL_free(ssl); SSL_CTX_free(ctx);
        return false;
    }

    *out_ctx = ctx;
    *out_ssl = ssl;
    return true;
}
\end{lstlisting}

\subsection{Branching I/O}

After connect:

\begin{lstlisting}[style=snippet,language=C]
SSL_CTX *ctx = NULL;
SSL *ssl = NULL;
if (u.https && !https_setup(s, u.host, &ctx, &ssl)) {
    close(s);
    return EXIT_FAILURE;
}

if (u.https)
    SSL_write(ssl, req, strlen(req));
else
    write(s, req, strlen(req));

while ((n = u.https
            ? SSL_read(ssl, buf, sizeof(buf))
            : read(s, buf, sizeof(buf))) > 0)
    fwrite(buf, 1, n, stdout);

if (u.https) {
    if (ssl) { SSL_shutdown(ssl); SSL_free(ssl); }
    if (ctx) SSL_CTX_free(ctx);
}
close(s);
\end{lstlisting}

Link with OpenSSL:

\begin{lstlisting}[style=snippet,language=bash]
gcc -Wall -Wextra -O2 -o urlget urlget.c -lssl -lcrypto
./urlget https://example.com/
\end{lstlisting}

\FullSourceStart{urlgetstep06}
\begin{lstlisting}[style=oldcode,language=C,firstnumber=1]
/* urlget.c -- Step 6: HTTPS via OpenSSL */
#include <stdio.h>
#include <stdlib.h>
#include <string.h>
#include <stdbool.h>
#include <stdint.h>
#include <unistd.h>
#include <sys/types.h>
#include <sys/socket.h>
#include <netdb.h>
\end{lstlisting}
\begin{lstlisting}[style=newcode,language=C,firstnumber=last]
#include <openssl/ssl.h>
#include <openssl/err.h>
\end{lstlisting}
\begin{lstlisting}[style=oldcode,language=C,firstnumber=last]

typedef struct {
    bool https;
    char host[256];
    uint16_t port;
    char path[512];
} url_t;

static inline int parse_url(const char *arg, url_t *u)
{
    char buf[1024];
    if (strlen(arg) >= sizeof(buf)) return -1;
    strcpy(buf, arg);

    char *p = buf;
    u->https = false;

    if (strncmp(p, "https://", 8) == 0) {
        u->https = true; p += 8;
    } else if (strncmp(p, "http://", 7) == 0) {
        p += 7;
    }

    char *sep = strpbrk(p, "/?");
    char *hostport = p;

    if (sep) {
        if (*sep == '?') {
            size_t ql = strlen(sep);
            if (ql + 2 > sizeof(u->path)) return -1;
            u->path[0] = '/';
            memcpy(u->path + 1, sep, ql + 1);
        } else {
            if (strlen(sep) >= sizeof(u->path)) return -1;
            strcpy(u->path, sep);
        }
        *sep = '\0';
    } else {
        strcpy(u->path, "/");
    }

    char *host = u->host;
    size_t hmax = sizeof(u->host) - 1;
    uint16_t def = u->https ? 443 : 80;
    u->port = def;

    if (*hostport == '[') {
        char *cl = strchr(hostport, ']');
        if (!cl) return -1;
        size_t hl = cl - (hostport + 1);
        if (hl >= hmax) return -1;
        memcpy(host, hostport + 1, hl);
        host[hl] = '\0';
        char *aft = cl + 1;
        if (*aft == ':')
            u->port = (uint16_t)strtol(aft + 1, NULL, 10);
    } else {
        char *col = strrchr(hostport, ':');
        if (col) {
            size_t hl = col - hostport;
            if (hl >= hmax) return -1;
            memcpy(host, hostport, hl);
            host[hl] = '\0';
            u->port = (uint16_t)strtol(col + 1, NULL, 10);
        } else {
            if (strlen(hostport) >= hmax) return -1;
            strcpy(host, hostport);
        }
    }

    if (u->port == 0) u->port = def;
    return 0;
}
\end{lstlisting}
\begin{lstlisting}[style=newcode,language=C,firstnumber=last]
static inline bool https_setup(int sock, const char *hostname,
                               SSL_CTX **out_ctx, SSL **out_ssl)
{
    SSL_CTX *ctx = SSL_CTX_new(TLS_client_method());
    if (!ctx) { ERR_print_errors_fp(stderr); return false; }

    SSL_CTX_set_verify(ctx, SSL_VERIFY_NONE, NULL);

    SSL *ssl = SSL_new(ctx);
    if (!ssl) { SSL_CTX_free(ctx); return false; }

    SSL_set_fd(ssl, sock);
    SSL_set_tlsext_host_name(ssl, hostname);

    if (SSL_connect(ssl) <= 0) {
        ERR_print_errors_fp(stderr);
        SSL_free(ssl); SSL_CTX_free(ctx);
        return false;
    }

    *out_ctx = ctx;
    *out_ssl = ssl;
    return true;
}
\end{lstlisting}
\begin{lstlisting}[style=oldcode,language=C,firstnumber=last]
static inline void usage(const char *progname)
{
    fprintf(stderr,
        "Usage: %s [options] <URL>\n\n"
        "Options:\n"
        "  -A, --user-agent STR     Set User-Agent\n"
        "  -T, --tor [proxy]        Use Tor (default 127.0.0.1:9050)\n"
        "  -H, --header STR         Add custom header (repeatable)\n"
        "  -d, --data DATA          POST data (@file supported)\n"
        "  -o, --output FILE        Write output to file\n"
        "  -u, --user USER:PASS     Basic authentication\n"
        "  -I, --head               HEAD request\n"
        "      --no-location        Do not follow redirects\n"
        "      --max-redirects N    Max redirects (default: 10)\n",
        progname);
    exit(EXIT_FAILURE);
}

int main(int argc, char *argv[])
{
    const char *user_agent =
        "Mozilla/5.0 (X11; Linux x86_64; rv:128.0) "
        "Gecko/20100101 Firefox/128.0";
    const char *default_tor = "127.0.0.1:9050";

    const char *url_arg = NULL;
    const char *tor_proxy = NULL;
    bool use_tor = false;
    bool follow_redirects = true;
    int max_redirects = 10;
    bool head_request = false;
    char extra_headers[4096] = {0};
    char post_data[8192] = {0};
    char output_file[256] = {0};
    char auth_header[256] = {0};

    for (int i = 1; i < argc; i++) {
        if (strcmp(argv[i], "-A") == 0 ||
            strcmp(argv[i], "--user-agent") == 0) {
            if (i + 1 < argc) user_agent = argv[++i];
            else usage(argv[0]);
        }
        else if (strcmp(argv[i], "-T") == 0 ||
                 strcmp(argv[i], "--tor") == 0) {
            use_tor = true;
            if (i + 1 < argc && argv[i + 1][0] != '-' &&
                strchr(argv[i + 1], ':')) {
                tor_proxy = argv[i + 1];
                i++;
            }
        }
        else if (strcmp(argv[i], "-H") == 0 ||
                 strcmp(argv[i], "--header") == 0) {
            if (i + 1 < argc) {
                strncat(extra_headers, argv[++i],
                        sizeof(extra_headers) -
                        strlen(extra_headers) - 3);
                strcat(extra_headers, "\r\n");
            } else usage(argv[0]);
        }
        else if (strcmp(argv[i], "-d") == 0 ||
                 strcmp(argv[i], "--data") == 0) {
            if (i + 1 < argc) {
                const char *val = argv[++i];
                if (val[0] == '@') {
                    FILE *f = fopen(val + 1, "r");
                    if (f) {
                        fread(post_data, 1,
                              sizeof(post_data) - 1, f);
                        fclose(f);
                    }
                } else {
                    strncpy(post_data, val,
                            sizeof(post_data) - 1);
                }
            } else usage(argv[0]);
        }
        else if (strcmp(argv[i], "-o") == 0 ||
                 strcmp(argv[i], "--output") == 0) {
            if (i + 1 < argc)
                strncpy(output_file, argv[++i],
                        sizeof(output_file) - 1);
            else usage(argv[0]);
        }
        else if (strcmp(argv[i], "-u") == 0 ||
                 strcmp(argv[i], "--user") == 0) {
            if (i + 1 < argc)
                snprintf(auth_header, sizeof(auth_header),
                    "Authorization: Basic %s\r\n", argv[++i]);
            else usage(argv[0]);
        }
        else if (strcmp(argv[i], "-I") == 0 ||
                 strcmp(argv[i], "--head") == 0) {
            head_request = true;
        }
        else if (strcmp(argv[i], "--no-location") == 0) {
            follow_redirects = false;
        }
        else if (strcmp(argv[i], "--max-redirects") == 0 &&
                 i + 1 < argc) {
            max_redirects = atoi(argv[++i]);
        }
        else if (url_arg == NULL) {
            url_arg = argv[i];
        }
    }

    if (url_arg == NULL)
        usage(argv[0]);
    if (use_tor && tor_proxy == NULL)
        tor_proxy = default_tor;

    url_t u = {0};
    if (parse_url(url_arg, &u) != 0) {
        fprintf(stderr, "Failed to parse URL\n");
        return EXIT_FAILURE;
    }

    struct addrinfo hints = {0};
    hints.ai_family   = AF_UNSPEC;
    hints.ai_socktype = SOCK_STREAM;
    char ps[6];
    snprintf(ps, sizeof(ps), "%u", u.port);

    struct addrinfo *res = NULL;
    if (getaddrinfo(u.host, ps, &hints, &res) != 0) {
        perror("getaddrinfo");
        return EXIT_FAILURE;
    }

    int s = -1;
    for (struct addrinfo *a = res; a; a = a->ai_next) {
        s = socket(a->ai_family, a->ai_socktype, a->ai_protocol);
        if (s < 0) continue;
        if (connect(s, a->ai_addr, a->ai_addrlen) == 0) break;
        close(s); s = -1;
    }
    freeaddrinfo(res);
    if (s < 0) { perror("connect"); return EXIT_FAILURE; }
\end{lstlisting}
\begin{lstlisting}[style=newcode,language=C,firstnumber=last]
    SSL_CTX *ctx = NULL;
    SSL *ssl = NULL;
    if (u.https && !https_setup(s, u.host, &ctx, &ssl)) {
        close(s);
        return EXIT_FAILURE;
    }

    char req[4096];
    snprintf(req, sizeof(req),
        "GET %s HTTP/1.1\r\n"
        "Host: %s\r\n"
        "User-Agent: %s\r\n"
        "Accept: */*\r\n"
        "Connection: close\r\n"
        "\r\n",
        u.path, u.host, user_agent);

    if (u.https)
        SSL_write(ssl, req, strlen(req));
    else
        write(s, req, strlen(req));

    char buf[8192];
    ssize_t n;
    while ((n = u.https
                ? SSL_read(ssl, buf, sizeof(buf))
                : read(s, buf, sizeof(buf))) > 0)
        fwrite(buf, 1, n, stdout);

    if (u.https) {
        if (ssl) { SSL_shutdown(ssl); SSL_free(ssl); }
        if (ctx) SSL_CTX_free(ctx);
    }
    close(s);

    (void)use_tor; (void)tor_proxy;
    (void)follow_redirects; (void)max_redirects;
    (void)head_request; (void)extra_headers;
    (void)post_data; (void)output_file; (void)auth_header;
\end{lstlisting}
\begin{lstlisting}[style=oldcode,language=C,firstnumber=last]
    return 0;
}
\end{lstlisting}
\FullSourceStop

\section{Step 7: Tor / SOCKS5 Support}

Direct connections reveal your IP to the destination. With \texttt{-T}, we
instead open a TCP connection to a local Tor SOCKS5 proxy (default
\texttt{127.0.0.1:9050}) and ask \emph{it} to reach the real host. The HTTP
or HTTPS session then rides on that tunnel.

\subsection{Why SOCKS5 (not HTTP CONNECT)?}

Tor's default application proxy speaks SOCKS5. Domain names can be sent to
the proxy (\texttt{ATYP = 0x03}), so DNS resolution also happens inside Tor
when configured that way---important for anonymity.

\subsection{Minimal Handshake}

\begin{enumerate}[leftmargin=1.4em]
    \item \textbf{Greeting:} version~5, one method ``no authentication''
          (\texttt{0x00}).
    \item \textbf{Server reply:} version~5, method \texttt{0x00} accepted.
    \item \textbf{Connect request:} \texttt{CMD=CONNECT}, address type domain,
          hostname length + bytes, port in network byte order.
    \item \textbf{Reply:} \texttt{REP=0x00} means success. (We ignore the
          bind address fields for a minimal client.)
\end{enumerate}

\subsection{New Code --- \texttt{socks5\_connect}}

\begin{lstlisting}[style=snippet,language=C]
static int socks5_connect(int sock, const char *host, uint16_t port)
{
    unsigned char buf[512];
    size_t hlen = strlen(host);
    if (hlen > 255) return -1;

    /* greeting: VER=5, NMETHODS=1, METHOD=0 (no auth) */
    buf[0] = 0x05; buf[1] = 0x01; buf[2] = 0x00;
    if (write(sock, buf, 3) != 3) return -1;

    if (read(sock, buf, 2) != 2 || buf[0] != 0x05 || buf[1] != 0x00)
        return -1;

    /* CONNECT request with domain name */
    buf[0] = 0x05; buf[1] = 0x01; buf[2] = 0x00; buf[3] = 0x03;
    buf[4] = (unsigned char)hlen;
    memcpy(buf + 5, host, hlen);
    buf[5 + hlen] = port >> 8;
    buf[6 + hlen] = port & 0xff;

    if (write(sock, buf, 7 + hlen) != (ssize_t)(7 + hlen)) return -1;
    if (read(sock, buf, 4) != 4 || buf[1] != 0x00) return -1;
    return 0;
}
\end{lstlisting}

\subsection{Wiring Into \texttt{main}}

When Tor is enabled, DNS/\texttt{connect} targets the \emph{proxy}, not the
origin. After the TCP connect succeeds, run the SOCKS5 handshake with the
real host and port from \texttt{url\_t}. TLS (if any) still uses the real
hostname for SNI.

\begin{lstlisting}[style=snippet,language=C]
const char *connect_host = use_tor ? tor_proxy : u.host;
uint16_t connect_port = use_tor ? 9050 : u.port;

if (use_tor) {
    char pbuf[256];
    strncpy(pbuf, tor_proxy, sizeof(pbuf) - 1);
    char *colon = strrchr(pbuf, ':');
    if (colon) {
        *colon = '\0';
        connect_port = (uint16_t)strtol(colon + 1, NULL, 10);
        connect_host = pbuf;
    }
}
/* getaddrinfo(connect_host) + connect ... */
if (use_tor && socks5_connect(s, u.host, u.port) != 0) {
    close(s);
    return EXIT_FAILURE;
}
\end{lstlisting}

\FullSourceStart{urlgetstep07}
\begin{lstlisting}[style=oldcode,language=C,firstnumber=1]
/* urlget.c -- Step 7: Tor / SOCKS5 */
#include <stdio.h>
#include <stdlib.h>
#include <string.h>
#include <stdbool.h>
#include <stdint.h>
#include <unistd.h>
#include <sys/types.h>
#include <sys/socket.h>
#include <netdb.h>
#include <openssl/ssl.h>
#include <openssl/err.h>

typedef struct {
    bool https;
    char host[256];
    uint16_t port;
    char path[512];
} url_t;

static inline int parse_url(const char *arg, url_t *u)
{
    char buf[1024];
    if (strlen(arg) >= sizeof(buf)) return -1;
    strcpy(buf, arg);
    char *p = buf;
    u->https = false;
    if (strncmp(p, "https://", 8) == 0) { u->https = true; p += 8; }
    else if (strncmp(p, "http://", 7) == 0) { p += 7; }
    char *sep = strpbrk(p, "/?");
    char *hostport = p;
    if (sep) {
        if (*sep == '?') {
            size_t ql = strlen(sep);
            if (ql + 2 > sizeof(u->path)) return -1;
            u->path[0] = '/';
            memcpy(u->path + 1, sep, ql + 1);
        } else {
            if (strlen(sep) >= sizeof(u->path)) return -1;
            strcpy(u->path, sep);
        }
        *sep = '\0';
    } else strcpy(u->path, "/");
    char *host = u->host;
    size_t hmax = sizeof(u->host) - 1;
    uint16_t def = u->https ? 443 : 80;
    u->port = def;
    if (*hostport == '[') {
        char *cl = strchr(hostport, ']');
        if (!cl) return -1;
        size_t hl = cl - (hostport + 1);
        if (hl >= hmax) return -1;
        memcpy(host, hostport + 1, hl);
        host[hl] = '\0';
        if (*(cl + 1) == ':')
            u->port = (uint16_t)strtol(cl + 2, NULL, 10);
    } else {
        char *col = strrchr(hostport, ':');
        if (col) {
            size_t hl = col - hostport;
            if (hl >= hmax) return -1;
            memcpy(host, hostport, hl);
            host[hl] = '\0';
            u->port = (uint16_t)strtol(col + 1, NULL, 10);
        } else {
            if (strlen(hostport) >= hmax) return -1;
            strcpy(host, hostport);
        }
    }
    if (u->port == 0) u->port = def;
    return 0;
}
\end{lstlisting}
\begin{lstlisting}[style=newcode,language=C,firstnumber=last]
static int socks5_connect(int sock, const char *host, uint16_t port)
{
    unsigned char buf[512];
    size_t hlen = strlen(host);
    if (hlen > 255) return -1;

    buf[0] = 0x05; buf[1] = 0x01; buf[2] = 0x00;
    if (write(sock, buf, 3) != 3) return -1;
    if (read(sock, buf, 2) != 2 || buf[0] != 0x05 || buf[1] != 0x00)
        return -1;

    buf[0] = 0x05; buf[1] = 0x01; buf[2] = 0x00; buf[3] = 0x03;
    buf[4] = (unsigned char)hlen;
    memcpy(buf + 5, host, hlen);
    buf[5 + hlen] = port >> 8;
    buf[6 + hlen] = port & 0xff;
    if (write(sock, buf, 7 + hlen) != (ssize_t)(7 + hlen)) return -1;
    if (read(sock, buf, 4) != 4 || buf[1] != 0x00) return -1;
    return 0;
}
\end{lstlisting}
\begin{lstlisting}[style=oldcode,language=C,firstnumber=last]
static inline bool https_setup(int sock, const char *hostname,
                               SSL_CTX **out_ctx, SSL **out_ssl)
{
    SSL_CTX *ctx = SSL_CTX_new(TLS_client_method());
    if (!ctx) { ERR_print_errors_fp(stderr); return false; }
    SSL_CTX_set_verify(ctx, SSL_VERIFY_NONE, NULL);
    SSL *ssl = SSL_new(ctx);
    if (!ssl) { SSL_CTX_free(ctx); return false; }
    SSL_set_fd(ssl, sock);
    SSL_set_tlsext_host_name(ssl, hostname);
    if (SSL_connect(ssl) <= 0) {
        ERR_print_errors_fp(stderr);
        SSL_free(ssl); SSL_CTX_free(ctx);
        return false;
    }
    *out_ctx = ctx;
    *out_ssl = ssl;
    return true;
}

static inline void usage(const char *progname)
{
    fprintf(stderr,
        "Usage: %s [options] <URL>\n\n"
        "Options:\n"
        "  -A, --user-agent STR     Set User-Agent\n"
        "  -T, --tor [proxy]        Use Tor (default 127.0.0.1:9050)\n"
        "  -H, --header STR         Add custom header (repeatable)\n"
        "  -d, --data DATA          POST data (@file supported)\n"
        "  -o, --output FILE        Write output to file\n"
        "  -u, --user USER:PASS     Basic authentication\n"
        "  -I, --head               HEAD request\n"
        "      --no-location        Do not follow redirects\n"
        "      --max-redirects N    Max redirects (default: 10)\n",
        progname);
    exit(EXIT_FAILURE);
}

int main(int argc, char *argv[])
{
    const char *user_agent =
        "Mozilla/5.0 (X11; Linux x86_64; rv:128.0) "
        "Gecko/20100101 Firefox/128.0";
    const char *default_tor = "127.0.0.1:9050";
    const char *url_arg = NULL;
    const char *tor_proxy = NULL;
    bool use_tor = false;
    bool follow_redirects = true;
    int max_redirects = 10;
    bool head_request = false;
    char extra_headers[4096] = {0};
    char post_data[8192] = {0};
    char output_file[256] = {0};
    char auth_header[256] = {0};

    for (int i = 1; i < argc; i++) {
        if (strcmp(argv[i], "-A") == 0 ||
            strcmp(argv[i], "--user-agent") == 0) {
            if (i + 1 < argc) user_agent = argv[++i];
            else usage(argv[0]);
        }
        else if (strcmp(argv[i], "-T") == 0 ||
                 strcmp(argv[i], "--tor") == 0) {
            use_tor = true;
            if (i + 1 < argc && argv[i + 1][0] != '-' &&
                strchr(argv[i + 1], ':')) {
                tor_proxy = argv[i + 1];
                i++;
            }
        }
        else if (strcmp(argv[i], "-H") == 0 ||
                 strcmp(argv[i], "--header") == 0) {
            if (i + 1 < argc) {
                strncat(extra_headers, argv[++i],
                        sizeof(extra_headers) -
                        strlen(extra_headers) - 3);
                strcat(extra_headers, "\r\n");
            } else usage(argv[0]);
        }
        else if (strcmp(argv[i], "-d") == 0 ||
                 strcmp(argv[i], "--data") == 0) {
            if (i + 1 < argc) {
                const char *val = argv[++i];
                if (val[0] == '@') {
                    FILE *f = fopen(val + 1, "r");
                    if (f) {
                        fread(post_data, 1,
                              sizeof(post_data) - 1, f);
                        fclose(f);
                    }
                } else
                    strncpy(post_data, val, sizeof(post_data) - 1);
            } else usage(argv[0]);
        }
        else if (strcmp(argv[i], "-o") == 0 ||
                 strcmp(argv[i], "--output") == 0) {
            if (i + 1 < argc)
                strncpy(output_file, argv[++i],
                        sizeof(output_file) - 1);
            else usage(argv[0]);
        }
        else if (strcmp(argv[i], "-u") == 0 ||
                 strcmp(argv[i], "--user") == 0) {
            if (i + 1 < argc)
                snprintf(auth_header, sizeof(auth_header),
                    "Authorization: Basic %s\r\n", argv[++i]);
            else usage(argv[0]);
        }
        else if (strcmp(argv[i], "-I") == 0 ||
                 strcmp(argv[i], "--head") == 0)
            head_request = true;
        else if (strcmp(argv[i], "--no-location") == 0)
            follow_redirects = false;
        else if (strcmp(argv[i], "--max-redirects") == 0 &&
                 i + 1 < argc)
            max_redirects = atoi(argv[++i]);
        else if (url_arg == NULL)
            url_arg = argv[i];
    }

    if (url_arg == NULL) usage(argv[0]);
    if (use_tor && tor_proxy == NULL) tor_proxy = default_tor;

    url_t u = {0};
    if (parse_url(url_arg, &u) != 0) {
        fprintf(stderr, "Failed to parse URL\n");
        return EXIT_FAILURE;
    }
\end{lstlisting}
\begin{lstlisting}[style=newcode,language=C,firstnumber=last]
    const char *connect_host = use_tor ? tor_proxy : u.host;
    uint16_t connect_port = use_tor ? 9050 : u.port;
    char pbuf[256];
    if (use_tor) {
        strncpy(pbuf, tor_proxy, sizeof(pbuf) - 1);
        char *colon = strrchr(pbuf, ':');
        if (colon) {
            *colon = '\0';
            connect_port = (uint16_t)strtol(colon + 1, NULL, 10);
            connect_host = pbuf;
        }
    }
\end{lstlisting}
\begin{lstlisting}[style=oldcode,language=C,firstnumber=last]
    struct addrinfo hints = {0};
    hints.ai_family   = AF_UNSPEC;
    hints.ai_socktype = SOCK_STREAM;
    char ps[6];
\end{lstlisting}
\begin{lstlisting}[style=newcode,language=C,firstnumber=last]
    snprintf(ps, sizeof(ps), "%u", connect_port);

    struct addrinfo *res = NULL;
    if (getaddrinfo(connect_host, ps, &hints, &res) != 0) {
\end{lstlisting}
\begin{lstlisting}[style=oldcode,language=C,firstnumber=last]
        perror("getaddrinfo");
        return EXIT_FAILURE;
    }

    int s = -1;
    for (struct addrinfo *a = res; a; a = a->ai_next) {
        s = socket(a->ai_family, a->ai_socktype, a->ai_protocol);
        if (s < 0) continue;
        if (connect(s, a->ai_addr, a->ai_addrlen) == 0) break;
        close(s); s = -1;
    }
    freeaddrinfo(res);
    if (s < 0) { perror("connect"); return EXIT_FAILURE; }
\end{lstlisting}
\begin{lstlisting}[style=newcode,language=C,firstnumber=last]
    if (use_tor && socks5_connect(s, u.host, u.port) != 0) {
        close(s);
        return EXIT_FAILURE;
    }
\end{lstlisting}
\begin{lstlisting}[style=oldcode,language=C,firstnumber=last]
    SSL_CTX *ctx = NULL;
    SSL *ssl = NULL;
    if (u.https && !https_setup(s, u.host, &ctx, &ssl)) {
        close(s);
        return EXIT_FAILURE;
    }

    char req[4096];
    snprintf(req, sizeof(req),
        "GET %s HTTP/1.1\r\n"
        "Host: %s\r\n"
        "User-Agent: %s\r\n"
        "Accept: */*\r\n"
        "Connection: close\r\n"
        "\r\n",
        u.path, u.host, user_agent);

    if (u.https)
        SSL_write(ssl, req, strlen(req));
    else
        write(s, req, strlen(req));

    char buf[8192];
    ssize_t n;
    while ((n = u.https
                ? SSL_read(ssl, buf, sizeof(buf))
                : read(s, buf, sizeof(buf))) > 0)
        fwrite(buf, 1, n, stdout);

    if (u.https) {
        if (ssl) { SSL_shutdown(ssl); SSL_free(ssl); }
        if (ctx) SSL_CTX_free(ctx);
    }
    close(s);
    (void)follow_redirects; (void)max_redirects;
    (void)head_request; (void)extra_headers;
    (void)post_data; (void)output_file; (void)auth_header;
    return 0;
}
\end{lstlisting}
\FullSourceStop

Requires a running Tor daemon (or any SOCKS5 proxy) on the configured
address:

\begin{lstlisting}[style=snippet,language=bash]
./urlget -T https://check.torproject.org/
./urlget -T 127.0.0.1:9050 https://example.com/
\end{lstlisting}

\include{step09_redirects}
\section{Step 9: Advanced Features}

Argument parsing already collected values for custom headers, POST bodies,
output files, basic auth, and HEAD requests. This step \emph{uses} those
values when building and sending the HTTP message, and when writing the
response.

\subsection{Method Selection}

\begin{lstlisting}[style=snippet,language=C]
const char *method = head_request ? "HEAD"
                     : (post_data[0] ? "POST" : "GET");
\end{lstlisting}

Priority: explicit \texttt{-I}/\texttt{--head} wins; otherwise any
\texttt{-d} data implies \texttt{POST}; default remains \texttt{GET}.

\subsection{Request Construction}

Optional header blocks are injected as pre-formatted strings ending in
\texttt{\textbackslash r\textbackslash n} (see the argument parser):

\begin{lstlisting}[style=snippet,language=C]
int len = snprintf(req, sizeof(req),
    "%s %s HTTP/1.1\r\n"
    "Host: %s\r\n"
    "User-Agent: %s\r\n"
    "%s"   /* extra_headers */
    "%s"   /* auth_header   */
    "Accept: */*\r\n"
    "Connection: close\r\n",
    method, u.path, u.host, user_agent,
    extra_headers, auth_header);

if (post_data[0]) {
    len += snprintf(req + len, sizeof(req) - len,
        "Content-Length: %zu\r\n\r\n%s",
        strlen(post_data), post_data);
} else {
    strcat(req + len, "\r\n");
}
\end{lstlisting}

\subsection{Output File (\texttt{-o})}

\begin{lstlisting}[style=snippet,language=C]
FILE *out = stdout;
if (output_file[0]) {
    out = fopen(output_file, "w");
    if (!out) { perror("fopen"); return EXIT_FAILURE; }
}
/* ... fwrite(buf, 1, n, out); ... */
if (out != stdout) fclose(out);
\end{lstlisting}

\subsection{Basic Authentication (\texttt{-u}) with Base64}

HTTP Basic auth expects
\texttt{Authorization: Basic <base64(user:pass)>}.
We accept curl-style \texttt{USER:PASS} and encode it ourselves with a small
RFC~4648 helper---no external process, no extra libraries.

\begin{lstlisting}[style=snippet,language=C]
static inline int base64_encode(const void *data, size_t len,
                                char *out, size_t out_sz)
{
    static const char T[] =
        "ABCDEFGHIJKLMNOPQRSTUVWXYZabcdefghijklmnopqrstuvwxyz"
        "0123456789+/";
    const unsigned char *s = data;
    size_t o = 0;

    for (size_t i = 0; i < len; i += 3) {
        unsigned v = (unsigned)s[i] << 16;
        int n = 1;
        if (i + 1 < len) { v |= (unsigned)s[i + 1] << 8; n = 2; }
        if (i + 2 < len) { v |= (unsigned)s[i + 2];      n = 3; }
        if (o + 4 >= out_sz) return -1;
        out[o++] = T[(v >> 18) & 63];
        out[o++] = T[(v >> 12) & 63];
        out[o++] = (n >= 2) ? T[(v >> 6) & 63] : '=';
        out[o++] = (n >= 3) ? T[v & 63]         : '=';
    }
    if (o >= out_sz) return -1;
    out[o] = '\0';
    return (int)o;
}
\end{lstlisting}

When parsing \texttt{-u} / \texttt{--user}:

\begin{lstlisting}[style=snippet,language=C]
char b64[200];
const char *cred = argv[++i];
if (base64_encode(cred, strlen(cred), b64, sizeof(b64)) < 0) {
    fprintf(stderr, "credentials too long\n");
    return EXIT_FAILURE;
}
snprintf(auth_header, sizeof(auth_header),
         "Authorization: Basic %s\r\n", b64);
\end{lstlisting}

The encoder walks the input three bytes at a time, emits four alphabet
characters, and writes \texttt{=} padding when the final group is short.
That is enough for credential strings and keeps the whole client
dependency-free beyond OpenSSL (which we already link for HTTPS).

\FullSourceStart{urlgetstep09}
\begin{lstlisting}[style=oldcode,language=C,firstnumber=1]
/* urlget.c -- Step 9: advanced request features */
#include <stdio.h>
#include <stdlib.h>
#include <string.h>
#include <stdbool.h>
#include <stdint.h>
#include <unistd.h>
#include <sys/types.h>
#include <sys/socket.h>
#include <netdb.h>
#include <openssl/ssl.h>
#include <openssl/err.h>
\end{lstlisting}
\begin{lstlisting}[style=newcode,language=C,firstnumber=last]
/* RFC 4648 Base64 for HTTP Basic auth (-u USER:PASS). */
static inline int base64_encode(const void *data, size_t len,
                                char *out, size_t out_sz)
{
    static const char T[] =
        "ABCDEFGHIJKLMNOPQRSTUVWXYZabcdefghijklmnopqrstuvwxyz"
        "0123456789+/";
    const unsigned char *s = data;
    size_t o = 0;

    for (size_t i = 0; i < len; i += 3) {
        unsigned v = (unsigned)s[i] << 16;
        int n = 1;
        if (i + 1 < len) { v |= (unsigned)s[i + 1] << 8; n = 2; }
        if (i + 2 < len) { v |= (unsigned)s[i + 2];      n = 3; }
        if (o + 4 >= out_sz) return -1;
        out[o++] = T[(v >> 18) & 63];
        out[o++] = T[(v >> 12) & 63];
        out[o++] = (n >= 2) ? T[(v >> 6) & 63] : '=';
        out[o++] = (n >= 3) ? T[v & 63]         : '=';
    }
    if (o >= out_sz) return -1;
    out[o] = '\0';
    return (int)o;
}
\end{lstlisting}
\begin{lstlisting}[style=oldcode,language=C,firstnumber=last]
typedef struct {
    bool https;
    char host[256];
    uint16_t port;
    char path[512];
} url_t;

static inline int parse_url(const char *arg, url_t *u)
{
    char buf[1024];
    if (strlen(arg) >= sizeof(buf)) return -1;
    strcpy(buf, arg);
    char *p = buf;
    u->https = false;
    if (strncmp(p, "https://", 8) == 0) { u->https = true; p += 8; }
    else if (strncmp(p, "http://", 7) == 0) { p += 7; }
    char *sep = strpbrk(p, "/?");
    char *hostport = p;
    if (sep) {
        if (*sep == '?') {
            size_t ql = strlen(sep);
            if (ql + 2 > sizeof(u->path)) return -1;
            u->path[0] = '/';
            memcpy(u->path + 1, sep, ql + 1);
        } else {
            if (strlen(sep) >= sizeof(u->path)) return -1;
            strcpy(u->path, sep);
        }
        *sep = '\0';
    } else strcpy(u->path, "/");
    char *host = u->host;
    size_t hmax = sizeof(u->host) - 1;
    uint16_t def = u->https ? 443 : 80;
    u->port = def;
    if (*hostport == '[') {
        char *cl = strchr(hostport, ']');
        if (!cl) return -1;
        size_t hl = cl - (hostport + 1);
        if (hl >= hmax) return -1;
        memcpy(host, hostport + 1, hl);
        host[hl] = '\0';
        if (*(cl + 1) == ':')
            u->port = (uint16_t)strtol(cl + 2, NULL, 10);
    } else {
        char *col = strrchr(hostport, ':');
        if (col) {
            size_t hl = col - hostport;
            if (hl >= hmax) return -1;
            memcpy(host, hostport, hl);
            host[hl] = '\0';
            u->port = (uint16_t)strtol(col + 1, NULL, 10);
        } else {
            if (strlen(hostport) >= hmax) return -1;
            strcpy(host, hostport);
        }
    }
    if (u->port == 0) u->port = def;
    return 0;
}

static int socks5_connect(int sock, const char *host, uint16_t port)
{
    unsigned char buf[512];
    size_t hlen = strlen(host);
    if (hlen > 255) return -1;
    buf[0] = 0x05; buf[1] = 0x01; buf[2] = 0x00;
    if (write(sock, buf, 3) != 3) return -1;
    if (read(sock, buf, 2) != 2 || buf[0] != 0x05 || buf[1] != 0x00)
        return -1;
    buf[0] = 0x05; buf[1] = 0x01; buf[2] = 0x00; buf[3] = 0x03;
    buf[4] = (unsigned char)hlen;
    memcpy(buf + 5, host, hlen);
    buf[5 + hlen] = port >> 8;
    buf[6 + hlen] = port & 0xff;
    if (write(sock, buf, 7 + hlen) != (ssize_t)(7 + hlen)) return -1;
    if (read(sock, buf, 4) != 4 || buf[1] != 0x00) return -1;
    return 0;
}

static inline bool https_setup(int sock, const char *hostname,
                               SSL_CTX **out_ctx, SSL **out_ssl)
{
    SSL_CTX *ctx = SSL_CTX_new(TLS_client_method());
    if (!ctx) { ERR_print_errors_fp(stderr); return false; }
    SSL_CTX_set_verify(ctx, SSL_VERIFY_NONE, NULL);
    SSL *ssl = SSL_new(ctx);
    if (!ssl) { SSL_CTX_free(ctx); return false; }
    SSL_set_fd(ssl, sock);
    SSL_set_tlsext_host_name(ssl, hostname);
    if (SSL_connect(ssl) <= 0) {
        ERR_print_errors_fp(stderr);
        SSL_free(ssl); SSL_CTX_free(ctx);
        return false;
    }
    *out_ctx = ctx;
    *out_ssl = ssl;
    return true;
}

static inline void usage(const char *progname)
{
    fprintf(stderr,
        "Usage: %s [options] <URL>\n\n"
        "Options:\n"
        "  -A, --user-agent STR     Set User-Agent\n"
        "  -T, --tor [proxy]        Use Tor (default 127.0.0.1:9050)\n"
        "  -H, --header STR         Add custom header (repeatable)\n"
        "  -d, --data DATA          POST data (@file supported)\n"
        "  -o, --output FILE        Write output to file\n"
        "  -u, --user USER:PASS     Basic authentication\n"
        "  -I, --head               HEAD request\n"
        "      --no-location        Do not follow redirects\n"
        "      --max-redirects N    Max redirects (default: 10)\n",
        progname);
    exit(EXIT_FAILURE);
}

int main(int argc, char *argv[])
{
    const char *user_agent =
        "Mozilla/5.0 (X11; Linux x86_64; rv:128.0) "
        "Gecko/20100101 Firefox/128.0";
    const char *default_tor = "127.0.0.1:9050";
    const char *url_arg = NULL;
    const char *tor_proxy = NULL;
    bool use_tor = false;
    bool follow_redirects = true;
    int max_redirects = 10;
    bool head_request = false;
    char extra_headers[4096] = {0};
    char post_data[8192] = {0};
    char output_file[256] = {0};
    char auth_header[256] = {0};

    for (int i = 1; i < argc; i++) {
        if (strcmp(argv[i], "-A") == 0 ||
            strcmp(argv[i], "--user-agent") == 0) {
            if (i + 1 < argc) user_agent = argv[++i];
            else usage(argv[0]);
        }
        else if (strcmp(argv[i], "-T") == 0 ||
                 strcmp(argv[i], "--tor") == 0) {
            use_tor = true;
            if (i + 1 < argc && argv[i + 1][0] != '-' &&
                strchr(argv[i + 1], ':')) {
                tor_proxy = argv[i + 1];
                i++;
            }
        }
        else if (strcmp(argv[i], "-H") == 0 ||
                 strcmp(argv[i], "--header") == 0) {
            if (i + 1 < argc) {
                strncat(extra_headers, argv[++i],
                        sizeof(extra_headers) -
                        strlen(extra_headers) - 3);
                strcat(extra_headers, "\r\n");
            } else usage(argv[0]);
        }
        else if (strcmp(argv[i], "-d") == 0 ||
                 strcmp(argv[i], "--data") == 0) {
            if (i + 1 < argc) {
                const char *val = argv[++i];
                if (val[0] == '@') {
                    FILE *f = fopen(val + 1, "r");
                    if (f) {
                        fread(post_data, 1,
                              sizeof(post_data) - 1, f);
                        fclose(f);
                    }
                } else
                    strncpy(post_data, val, sizeof(post_data) - 1);
            } else usage(argv[0]);
        }
        else if (strcmp(argv[i], "-o") == 0 ||
                 strcmp(argv[i], "--output") == 0) {
            if (i + 1 < argc)
                strncpy(output_file, argv[++i],
                        sizeof(output_file) - 1);
            else usage(argv[0]);
        }
\end{lstlisting}
\begin{lstlisting}[style=newcode,language=C,firstnumber=last]
        else if (strcmp(argv[i], "-u") == 0 ||
                 strcmp(argv[i], "--user") == 0) {
            if (i + 1 < argc) {
                char b64[200];
                const char *cred = argv[++i];
                if (base64_encode(cred, strlen(cred), b64,
                                  sizeof(b64)) < 0) {
                    fprintf(stderr, "credentials too long\n");
                    return EXIT_FAILURE;
                }
                snprintf(auth_header, sizeof(auth_header),
                         "Authorization: Basic %s\r\n", b64);
            } else usage(argv[0]);
        }
\end{lstlisting}
\begin{lstlisting}[style=oldcode,language=C,firstnumber=last]
        else if (strcmp(argv[i], "-I") == 0 ||
                 strcmp(argv[i], "--head") == 0)
            head_request = true;
        else if (strcmp(argv[i], "--no-location") == 0)
            follow_redirects = false;
        else if (strcmp(argv[i], "--max-redirects") == 0 &&
                 i + 1 < argc)
            max_redirects = atoi(argv[++i]);
        else if (url_arg == NULL)
            url_arg = argv[i];
    }

    if (url_arg == NULL) usage(argv[0]);
    if (use_tor && tor_proxy == NULL) tor_proxy = default_tor;
\end{lstlisting}
\begin{lstlisting}[style=newcode,language=C,firstnumber=last]
    FILE *out = stdout;
    if (output_file[0]) {
        out = fopen(output_file, "w");
        if (!out) { perror("fopen"); return EXIT_FAILURE; }
    }

    int redirect_count = 0;
    char current_url[1024];
    strncpy(current_url, url_arg, sizeof(current_url) - 1);

    while (redirect_count <= max_redirects) {
        url_t u = {0};
        if (parse_url(current_url, &u) != 0) {
            fprintf(stderr, "Failed to parse URL\n");
            return EXIT_FAILURE;
        }

        const char *connect_host = use_tor ? tor_proxy : u.host;
        uint16_t connect_port = use_tor ? 9050 : u.port;
        char pbuf[256];
        if (use_tor) {
            strncpy(pbuf, tor_proxy, sizeof(pbuf) - 1);
            char *colon = strrchr(pbuf, ':');
            if (colon) {
                *colon = '\0';
                connect_port =
                    (uint16_t)strtol(colon + 1, NULL, 10);
                connect_host = pbuf;
            }
        }

        struct addrinfo hints = {0};
        hints.ai_family   = AF_UNSPEC;
        hints.ai_socktype = SOCK_STREAM;
        char ps[6];
        snprintf(ps, sizeof(ps), "%u", connect_port);

        struct addrinfo *res = NULL;
        if (getaddrinfo(connect_host, ps, &hints, &res) != 0) {
            perror("getaddrinfo");
            return EXIT_FAILURE;
        }

        int s = -1;
        for (struct addrinfo *a = res; a; a = a->ai_next) {
            s = socket(a->ai_family, a->ai_socktype, a->ai_protocol);
            if (s < 0) continue;
            if (connect(s, a->ai_addr, a->ai_addrlen) == 0) break;
            close(s); s = -1;
        }
        freeaddrinfo(res);
        if (s < 0) { perror("connect"); return EXIT_FAILURE; }

        if (use_tor && socks5_connect(s, u.host, u.port) != 0) {
            close(s);
            return EXIT_FAILURE;
        }

        SSL_CTX *ctx = NULL;
        SSL *ssl = NULL;
        if (u.https && !https_setup(s, u.host, &ctx, &ssl)) {
            close(s);
            return EXIT_FAILURE;
        }

        const char *method = head_request ? "HEAD"
                             : (post_data[0] ? "POST" : "GET");

        char req[4096];
        int len = snprintf(req, sizeof(req),
            "%s %s HTTP/1.1\r\n"
            "Host: %s\r\n"
            "User-Agent: %s\r\n"
            "%s"
            "%s"
            "Accept: */*\r\n"
            "Connection: close\r\n",
            method, u.path, u.host, user_agent,
            extra_headers, auth_header);

        if (post_data[0]) {
            len += snprintf(req + len, sizeof(req) - len,
                "Content-Length: %zu\r\n\r\n%s",
                strlen(post_data), post_data);
        } else {
            strcat(req + len, "\r\n");
        }

        if (u.https)
            SSL_write(ssl, req, strlen(req));
        else
            write(s, req, strlen(req));

        char buf[8192];
        ssize_t n;
        while ((n = u.https
                    ? SSL_read(ssl, buf, sizeof(buf))
                    : read(s, buf, sizeof(buf))) > 0)
            fwrite(buf, 1, n, out);

        if (u.https) {
            if (ssl) { SSL_shutdown(ssl); SSL_free(ssl); }
            if (ctx) SSL_CTX_free(ctx);
        }
        close(s);

        if (follow_redirects && strstr(buf, "Location:") &&
            redirect_count < max_redirects) {
            char location[1024] = {0};
            char *loc = strstr(buf, "Location:");
            if (loc) sscanf(loc, "Location: %1023s", location);
            if (location[0]) {
                redirect_count++;
                strncpy(current_url, location,
                        sizeof(current_url) - 1);
                continue;
            }
        }
        break;
    }

    if (out != stdout) fclose(out);
\end{lstlisting}
\begin{lstlisting}[style=oldcode,language=C,firstnumber=last]
    return 0;
}
\end{lstlisting}
\FullSourceStop

\begin{lstlisting}[style=snippet,language=bash]
./urlget -I https://example.com/
./urlget -H "Accept-Language: en" https://example.com/
./urlget -d "name=Ada" -o resp.txt https://httpbin.org/post
# curl-style USER:PASS --- encoded inside urlget
./urlget -u 'user:pass' https://httpbin.org/basic-auth/user/pass
\end{lstlisting}

\include{step11_final_polish}

\end{document}
